\documentclass[11pt,twocolumn,a4paper]{article}

\usepackage[utf8]{inputenc}
\usepackage[T1]{fontenc}
\usepackage{lmodern}
\usepackage[margin=1in]{geometry}
\usepackage{amsmath, amssymb, amsthm, mathtools}
\usepackage{thmtools}
\usepackage{enumitem}
\usepackage{microtype}
\usepackage{booktabs}
\usepackage{graphicx}
\usepackage{caption}
\usepackage{subcaption}
\usepackage{hyperref}
\usepackage{cleveref}
\usepackage{xcolor}
\usepackage{tikz}
\usetikzlibrary{arrows.meta,positioning,calc}

\hypersetup{
  colorlinks=true, linkcolor=blue!50!black,
  citecolor=green!50!black, urlcolor=blue!60!black
}

\declaretheoremstyle[
  spaceabove=6pt, spacebelow=6pt,
  headfont=\bfseries, notefont=\mdseries, notebraces={(}{)},
  bodyfont=\itshape, postheadspace=1em
]{thmstyle}
\declaretheoremstyle[
  spaceabove=6pt, spacebelow=6pt,
  headfont=\bfseries, notefont=\mdseries, notebraces={(}{)},
  bodyfont=\normalfont, postheadspace=1em
]{defstyle}

\declaretheorem[style=thmstyle, name=Theorem, numberwithin=section]{theorem}
\declaretheorem[style=thmstyle, name=Proposition, sibling=theorem]{proposition}
\declaretheorem[style=thmstyle, name=Lemma, sibling=theorem]{lemma}
\declaretheorem[style=thmstyle, name=Corollary, sibling=theorem]{corollary}
\declaretheorem[style=defstyle, name=Definition, sibling=theorem]{definition}
\declaretheorem[style=defstyle, name=Axiom, sibling=theorem]{axiom}
\declaretheorem[style=defstyle, name=Remark, sibling=theorem]{remark}
\declaretheorem[style=defstyle, name=Example, sibling=theorem]{example}
\declaretheorem[style=defstyle, name=Principle, sibling=theorem]{principle}
\declaretheorem[style=defstyle, name=Construction, sibling=theorem]{construction}

% cleveref names for custom theorem-like environments (fixes Cref ranges)
\crefname{axiom}{Axiom}{Axioms}
\Crefname{axiom}{Axiom}{Axioms}
\crefname{proposition}{Proposition}{Propositions}
\Crefname{proposition}{Proposition}{Propositions}
\crefname{definition}{Definition}{Definitions}
\Crefname{definition}{Definition}{Definitions}
\crefname{remark}{Remark}{Remarks}
\Crefname{remark}{Remark}{Remarks}
\crefname{corollary}{Corollary}{Corollaries}
\Crefname{corollary}{Corollary}{Corollaries}

\newcommand{\R}{\mathbb{R}}
\newcommand{\N}{\mathbb{N}}
\newcommand{\K}{\mathcal{K}}
\newcommand{\D}{\mathcal{D}}
\newcommand{\M}{\mathcal{M}}
\newcommand{\X}{\mathcal{X}}
\newcommand{\Cell}{\mathcal{C}}
\newcommand{\Aper}{\mathcal{A}}
\newcommand{\Mode}{\mathfrak{M}}
\newcommand{\Method}{\mathfrak{P}}
\newcommand{\Goal}{\mathfrak{G}}
\newcommand{\receiver}{\mathsf{R}}
\newcommand{\agent}{\mathsf{A}}
\newcommand{\Sfloor}{S_{\flat}}
\newcommand{\eps}{\varepsilon}
\newcommand{\norm}[1]{\lVert#1\rVert}
\newcommand{\abs}[1]{\lvert#1\rvert}
\newcommand{\dd}{\,\mathrm{d}}
\newcommand{\Compose}{\diamond}
\newcommand{\Sk}{S_k}
\newcommand{\St}{S_t}
\newcommand{\Se}{S_e}
\newcommand{\Svec}{\mathbf{S}}
\newcommand{\qvec}{\mathbf{q}}
\newcommand{\Lag}{\mathcal{L}}
\newcommand{\Ham}{\mathcal{H}}
\newcommand{\Pot}{\Phi}
\newcommand{\kB}{k_{\mathrm{B}}}
\newcommand{\Kc}{K_c}
\newcommand{\tauP}{\tau_p}
\newcommand{\Ord}{R}
\newcommand{\Tc}{T_c}

\title{\textbf{Coordination Regimes of Synchronised Agents:\\
A Self-Contained Variational Calculus on the\\
S-Entropy Manifold, with External Validation}}

\author{Kundai Farai Sachikonye\\
Technical University of Munich\\
AIMe Registry for Artificial Intelligence\\
Maximus-von-Imhof-Forum 3, 85354 Freising, Germany\\
\texttt{kundai.sachikonye@wzw.tum.de}}

\date{\today}

\begin{document}
\maketitle

\begin{abstract}
We develop a self-contained variational calculus for the coordination of
synchronised agents. An agent is modelled as an overdamped (Onsager--Machlup)
dynamical system on a five-dimensional manifold
$\M=[0,1]\times[0,2\pi^2]\times[0,1]^3$ whose coordinates are an internal
Kuramoto order parameter $\Ord$, a phase variance $\sigma^2$, and three
\emph{S-entropy} coordinates $(\Sk,\St,\Se)$. Crucially, and in contrast to a
preliminary version of this work, every construct used here---the S-entropy
coordinates, the partition potential, the categorical aperture, and the
governing Lagrangian---is \emph{derived within the paper} from three stated
axioms (bounded phase space, categorical exclusion, finite resolution); the
manuscript depends on no external or unpublished result. We prove: (i) a
triple-equivalence identity establishing the S-entropy coordinates from
maximum-entropy counting; (ii) the explicit Euler--Lagrange equations of the
agent action, given in full rather than asserted; (iii) the Kuramoto
synchronisation bifurcation, with the \emph{corrected} critical coupling
$\Kc=2/[\pi g(0)]$ (for a Gaussian frequency law, $\Kc=\sqrt{2\pi}\,\sigma_\omega
\,(2/\pi)\approx 1.596\,\sigma_\omega$), replacing an erroneous coefficient in
the preliminary version; (iv) a formal, analogy-free definition of
\emph{partition extinction} and an observable-commutation theorem that
separates functional from ontological indistinguishability; and (v) a
coordination calculus for ensembles. We are explicit that the five coordination
\emph{bands} on $\Ord$ are operational classification cut-offs, not
phase transitions: the only genuine transition in the model is the Kuramoto
bifurcation. Rather than the self-referential numerical checks of the
preliminary version, we report \emph{nine external experiments} against public
data (atomic shell structure, the Kuramoto onset by direct simulation, the
Wiedemann--Franz constant across eleven metals, sleep-EEG order parameters from
PhysioNet, enzyme kinetics, antidepressant meta-analysis, superconductor
transport, and representation-disjoint classifiers on a benchmark dataset).
Two of these external tests forced corrections to the theory and are reported
as such. A Scope and Limitations section delimits the defensible claims.
\end{abstract}

\tableofcontents


%==========================================================
\part{Foundations}
%==========================================================

\section{Introduction}
\label{sec:intro}

\subsection{Aim and the independence requirement}

This paper develops a variational calculus for the coordination of
synchronised agents and tests it against external data. A preliminary version
of this work was written as one entry in a larger programme and, in
consequence, relied on constructions (an ``S-entropy'' coordinate system, a
``cell-truth'' receiver calculus, a partition Lagrangian) that were defined in
companion manuscripts rather than in the paper itself. Reviewers correctly
identified this as a defect: a reader cannot verify a theory whose primitives
live elsewhere. The present version is rewritten to be \emph{mathematically
self-contained}. Every object is constructed here from three axioms; every
theorem is proved from those axioms and standard mathematics; and no claim
depends on an unpublished result.

A second defect of the preliminary version was its empirical section, which
reported fifty numerical ``experiments'' agreeing with theory ``to machine
precision.'' As reviewers noted, those checks evaluated closed-form predictions
against the same closed forms; agreement to $10^{-16}$ reflects floating-point
rounding of identities, not validation. We replace that section with nine
\emph{external} experiments (\Cref{sec:external}) that compare the theory to
data it did not generate. Two of the nine contradicted the preliminary theory
and forced corrections, which we make explicit. We regard the presence of
genuine corrections as evidence that the remaining confirmations are
meaningful.

\subsection{The coordination phenomenon}

Agents with markedly different internal organisation routinely converge on
common outcomes: drivers avoid collisions, independent laboratories replicate a
finding through different instruments, members of an ensemble lock to a shared
tempo. The classical account grounds such coordination in common
knowledge or focal points~\cite{lewis1969,schelling1960,aumann1976}, an
idealisation that is unreachable for agents with mutually inaccessible
representations~\cite{fagin1995}. We pursue a different account: coordination is
synchronisation of a collective order parameter, and the deepest coordinated
state is one in which inter-agent distinctions cease to be operative. We make
this precise and, where the data permit, test it.

\subsection{Contributions, and what is and is not new}

We separate genuinely new contributions from imported standard results, as
reviewers requested.

\emph{Imported (standard, used as tools):} the Kuramoto model and its
mean-field bifurcation~\cite{kuramoto1975,strogatz2000,acebron2005,dorfler2014};
the Onsager--Machlup action for overdamped dynamics~\cite{onsager1953,
machlup1953,graham1977}; the Banach fixed-point theorem~\cite{banach1922};
maximum-entropy inference~\cite{shannon1948,coverthomas2006}; and the
Wiedemann--Franz law and partition-function counting from standard physics.

\emph{New here:} (1) the derivation of the S-entropy coordinate manifold from
three axioms via a triple-equivalence identity (\Cref{sec:axioms}); (2) the
explicit partition potential and the full Euler--Lagrange system of the agent
action (\Cref{sec:dynamics}); (3) a precise, analogy-free definition of
partition extinction together with an observable-commutation theorem
distinguishing functional from ontological indistinguishability
(\Cref{sec:extinction}); (4) the identification of the coordinated/phase-locked
regime with partition extinction \emph{as a structural correspondence, not an
identity}; and (5) the external-validation programme of \Cref{sec:external},
including the corrected critical coupling.

\subsection{A note on the five ``regimes''}

The preliminary version asserted that the order parameter $\Ord$ separates five
regimes at $\Ord\in\{0.3,0.5,0.8,0.95\}$ and that these boundaries are
second-order phase transitions. We retract that assertion. The boundary values
are \emph{operational classification cut-offs}---convenient labels for bands of
a continuous variable---and we show below that they are not phase transitions:
the partition potential has no curvature sign-change at them. The only genuine
transition in the model is the Kuramoto synchronisation bifurcation at $\Kc$
(\Cref{sec:dynamics}). This honest demotion removes a circular argument present
in the preliminary version (regimes defined by thresholds, thresholds then
``derived'' from the regimes) and is consistent with the fact that companion
analyses used different numerical cut-offs for the same qualitative bands.

\subsection{Organisation}

\textbf{Part I} states the axioms and derives the S-entropy manifold, the
partition potential, the agent Lagrangian, and the explicit Euler--Lagrange
equations. \textbf{Part II} develops single-agent dynamics and the Kuramoto
bifurcation with the corrected critical coupling. \textbf{Part III} develops
ensembles and the coordination calculus. \textbf{Part IV} gives the
analogy-free theory of partition extinction. \textbf{Part V} reports the nine
external experiments, the Scope and Limitations, and the conclusion.

%==========================================================
\section{Axioms and the S-entropy manifold}
\label{sec:axioms}
%==========================================================

We adopt three axioms. They are deliberately weak; their role is to fix the
state space and its counting structure, not to encode the conclusions.

\begin{axiom}[Bounded phase space]\label{ax:bounded}
The accessible phase space $\Omega$ of an agent has finite measure,
$\mu(\Omega)<\infty$.
\end{axiom}

\begin{axiom}[Categorical exclusion]\label{ax:null}
At each instant the agent occupies exactly one of a finite set $\mathcal{C}$ of
mutually exclusive categorical states; there is no null state.
\end{axiom}

\begin{axiom}[Finite resolution]\label{ax:resolution}
Every observation partitions $\Omega$ into finitely many distinguishable cells;
states closer than a resolution $\delta>0$ are indistinguishable.
\end{axiom}

\begin{remark}
\Cref{ax:bounded,ax:null,ax:resolution} are satisfied by essentially every
physical or computational system, which is both their strength (generality) and
their weakness (they do not, by themselves, discriminate among systems). We
flag this explicitly; the discriminating content of the theory is in the
dynamics built on top of them, and in the external tests of
\Cref{sec:external}, not in the axioms.
\end{remark}

\subsection{Partition coordinates and capacity}

\begin{definition}[Partition coordinates]\label{def:partition-coords}
Under \Cref{ax:bounded,ax:null,ax:resolution} a bounded three-dimensional
oscillatory state is labelled by a tuple $(n,\ell,m,s)$ with $n\in\N$,
$0\le\ell<n$, $-\ell\le m\le\ell$, and $s\in\{-\tfrac12,+\tfrac12\}$, arising
from the nested closure conditions of bounded oscillation (radial depth $n$,
angular index $\ell$, orientation $m$, chirality $s$).
\end{definition}

\begin{theorem}[Shell capacity]\label{thm:capacity}
The number of partition states at depth $n$ is $C(n)=2n^2$.
\end{theorem}

\begin{proof}
For each $\ell\in\{0,\dots,n-1\}$ there are $2\ell+1$ values of $m$ and two
values of $s$, so $C(n)=2\sum_{\ell=0}^{n-1}(2\ell+1)=2n^2$.
\end{proof}

\begin{remark}
We test $C(n)=2n^2$ against the real periodic table in \Cref{sec:external}
(EXT01). The formula is exact as a level degeneracy; we will be careful there
to distinguish it from the question of shell \emph{closure}, which requires the
Aufbau filling order and which $2n^2$ alone does not determine.
\end{remark}

\subsection{The triple equivalence and the S-entropy coordinates}

\begin{definition}[Three entropies]\label{def:three-entropies}
For $M$ independent modes each occupying one of $n$ distinguishable states:
the \emph{oscillatory} entropy $S_{\mathrm{osc}}=-\kB\sum_j p_j\ln p_j$ over
mode occupations $p_j$; the \emph{categorical} entropy $S_{\mathrm{cat}}=
\kB\ln W$ over the number $W$ of consistent categorical microstates; and the
\emph{partition} entropy $S_{\mathrm{part}}=\kB\ln C(n)^M$.
\end{definition}

\begin{theorem}[Triple equivalence]\label{thm:triple}
At maximum entropy over the $C(n)$ states of depth $n$, the three entropies of
\Cref{def:three-entropies} coincide:
\[
S_{\mathrm{osc}}=S_{\mathrm{cat}}=S_{\mathrm{part}}=\kB M\ln n+\text{const},
\]
where the constant absorbs the $n$-independent term $\kB M\ln 2$.
\end{theorem}

\begin{proof}
By the maximum-entropy principle~\cite{shannon1948,coverthomas2006}, the
distribution over the $C(n)=2n^2$ states maximising $S_{\mathrm{osc}}$ subject
only to normalisation is uniform, $p_j=1/C(n)$, giving
$S_{\mathrm{osc}}=\kB\ln C(n)$. For $M$ independent modes,
$S_{\mathrm{osc}}^{(M)}=\kB M\ln C(n)$. Boltzmann counting gives
$W=C(n)^M$, so $S_{\mathrm{cat}}=\kB\ln W=\kB M\ln C(n)$. By definition
$S_{\mathrm{part}}=\kB M\ln C(n)$. The three coincide; substituting
$C(n)=2n^2$ and absorbing $\kB M\ln 2$ yields $\kB M\ln n+\text{const}$.
\end{proof}

\begin{remark}
\Cref{thm:triple} is the internal justification for treating ``oscillatory,''
``categorical,'' and ``partition'' descriptions interchangeably: they are three
counts of the same configuration set.
\end{remark}

\begin{definition}[S-entropy coordinates]\label{def:s-entropy}
Normalising the entropy contributions from configuration, timing, and energy by
the maximal attainable entropy $S_{\max}=\kB M\ln n_{\max}$ gives the
\emph{S-entropy coordinates}
\[
\Svec=(\Sk,\St,\Se)\in[0,1]^3,
\]
where $\Sk$ (knowledge/configuration), $\St$ (temporal), and $\Se$ (energetic/
evolution) are the three normalised components. The unit cube $[0,1]^3$ is the
\emph{S-entropy cube}.
\end{definition}

\begin{remark}[What the S-entropy coordinates are, and are not]
\label{rem:s-content}
The S-entropy coordinates measure the \emph{shape} of the agent's distribution
over its partition states---they are normalised entropies, not content-laden,
ecologically interpreted variables. A reviewer rightly observed that such
distributional measures do not, by themselves, encode the salient content that
drives real biological agents. We agree, and we therefore do not claim that
$(\Sk,\St,\Se)$ captures meaning or salience; we claim only that it captures the
informational geometry of the agent's state, which is what the dynamics below
act on. The empirical reach of the framework is correspondingly limited
(\Cref{sec:scope}).
\end{remark}

\subsection{The state manifold}

\begin{definition}[State manifold]\label{def:manifold}
The agent state manifold is
\[
\M=[0,1]\times[0,2\pi^2]\times[0,1]^3,\qquad
\qvec=(\Ord,\sigma^2,\Sk,\St,\Se),
\]
with $\Ord\in[0,1]$ the internal Kuramoto order parameter, $\sigma^2\in
[0,2\pi^2]$ the phase variance, and $(\Sk,\St,\Se)$ the S-entropy coordinates.
$\M$ is compact (a product of compact intervals), smooth in its interior,
connected, and a manifold with corners; we equip it with the block-diagonal
metric $g=\operatorname{diag}(\mu_R,\mu_\sigma,1,1,1)$, where $\mu_R,\mu_\sigma>0$
are inertial parameters and the S-entropy block carries the Euclidean metric.
\end{definition}

\begin{remark}
The precise topological description requested by reviewers: $\M$ is compact and
connected with non-empty boundary (corners), so it is not boundaryless; the
dynamics below are gradient flows on its interior with reflecting behaviour at
the boundary enforced by the boundary potential $V_{\mathrm{ent}}$
(\Cref{def:potential}).
\end{remark}

%==========================================================
\section{Single-agent dynamics: potential, action, and explicit
Euler--Lagrange equations}
\label{sec:dynamics}
%==========================================================

\subsection{The partition potential}

\begin{definition}[Partition potential]\label{def:potential}
The agent's dynamics are generated by the scalar potential
$\Pot:\M\to\R$,
\[
\Pot(\qvec)=V_{\mathrm{sync}}+V_{\mathrm{var}}+V_{\mathrm{SF}}+V_{\mathrm{ent}},
\]
with the four terms
\begin{align}
V_{\mathrm{sync}}(\Ord;K)&=\tfrac12(\Kc-K)\Ord^2+\tfrac14 K\Ord^4,\\
V_{\mathrm{var}}(\sigma^2)&=\kB T\,\sigma^2+\frac{\kB T}{K\sigma^2},\\
V_{\mathrm{SF}}(\Ord,\sigma^2)&=-\alpha\,\Ord\,e^{-\sigma^2/2\pi^2},\\
V_{\mathrm{ent}}(\Svec)&=\frac{\lambda}{2}\!\sum_{i\in\{k,t,e\}}\!\!\big[
(\max(0,-S_i))^2+(\max(0,S_i-1))^2\big]-T_{\mathrm{eff}}H(\Svec),
\end{align}
where $K$ is the intra-agent coupling, $\Kc$ the critical coupling
(\Cref{thm:Kc}), $\alpha>0$ a structural coupling, $\lambda\gg1$ a wall
stiffness enforcing $\Svec\in[0,1]^3$, and $H(\Svec)=-\sum_i[S_i\ln S_i+
(1-S_i)\ln(1-S_i)]$ the binary entropy.
\end{definition}

\begin{remark}[Provenance of each term]
$V_{\mathrm{sync}}$ is the Landau form whose minimisation reproduces the
Kuramoto order parameter (\Cref{thm:bifurcation}); $V_{\mathrm{var}}$ is a
convex well in $\sigma^2$ producing a variance floor; $V_{\mathrm{SF}}$ couples
coherence to low variance; $V_{\mathrm{ent}}$ confines the S-entropy
coordinates to the cube and drives them, absent forcing, toward the
maximum-entropy centre $(\tfrac12,\tfrac12,\tfrac12)$. No term encodes the
regime cut-offs; the bands are read off $\Ord$ afterward and play no role in
$\Pot$.
\end{remark}

\subsection{The Onsager--Machlup action}

The Kuramoto and relaxational dynamics relevant here are first-order
(overdamped), not conservative; the appropriate variational principle is
therefore the Onsager--Machlup action, whose extremum is the most-probable path
of the overdamped Langevin dynamics~\cite{onsager1953,machlup1953,graham1977}.

\begin{definition}[Overdamped dynamics]\label{def:langevin}
The agent obeys $\dot q^i=-g^{ij}\partial_j\Pot+\eta^i(t)$ with Gaussian white
noise $\langle\eta^i(t)\eta^j(t')\rangle=2D\,g^{ij}\delta(t-t')$.
\end{definition}

\begin{definition}[Agent action]\label{def:action}
The Onsager--Machlup Lagrangian and action are
\[
\Lag(\qvec,\dot\qvec)=\frac{1}{4D}\,g_{ij}\big(\dot q^i+g^{ik}\partial_k\Pot\big)
\big(\dot q^j+g^{jl}\partial_l\Pot\big)-\tfrac12\nabla^2\Pot,\quad
\mathcal{S}[\qvec]=\!\int_0^T\!\Lag\dd t,
\]
where $\nabla^2\Pot=g^{ij}\partial_i\partial_j\Pot$.
\end{definition}

\subsection{Explicit Euler--Lagrange derivation}

The preliminary version asserted the equations of motion as ``standard calculus
of variations.'' We give the derivation in full.

\begin{theorem}[Euler--Lagrange equations of the agent action]\label{thm:EL}
The stationarity condition $\delta\mathcal{S}/\delta q^i=0$ for the action of
\Cref{def:action} yields, for each coordinate $i$,
\begin{equation}
\frac{1}{2D}g_{ij}\ddot q^j
=\frac{1}{2D}\,\partial_i\!\Big(\tfrac12 g_{kl}\,b^k b^l\Big)
-\tfrac12\,\partial_i\nabla^2\Pot,
\qquad b^k:=g^{km}\partial_m\Pot,
\label{eq:EL-second-order}
\end{equation}
and the deterministic most-probable path in the zero-noise (gradient) limit is
\begin{equation}
\dot q^i=-g^{ij}\partial_j\Pot.
\label{eq:gradflow}
\end{equation}
\end{theorem}

\begin{proof}
Write $f^i:=\dot q^i+b^i$ with $b^i=g^{ik}\partial_k\Pot$, so
$\Lag=\frac{1}{4D}g_{ij}f^if^j-\tfrac12\nabla^2\Pot$. The Euler--Lagrange
operator is $\frac{d}{dt}\frac{\partial\Lag}{\partial\dot q^i}
-\frac{\partial\Lag}{\partial q^i}$.

\emph{Momentum term.} Since $\partial f^j/\partial\dot q^i=\delta^j_i$,
\[
\frac{\partial\Lag}{\partial\dot q^i}=\frac{1}{2D}g_{ij}f^j
=\frac{1}{2D}g_{ij}(\dot q^j+b^j),
\]
hence
\[
\frac{d}{dt}\frac{\partial\Lag}{\partial\dot q^i}
=\frac{1}{2D}g_{ij}\big(\ddot q^j+\dot b^j\big),
\qquad \dot b^j=\partial_m b^j\,\dot q^m .
\]

\emph{Coordinate term.} Using $\partial f^j/\partial q^i=\partial_i b^j$,
\[
\frac{\partial\Lag}{\partial q^i}
=\frac{1}{2D}g_{kl}f^k\,\partial_i b^l
-\tfrac12\,\partial_i\nabla^2\Pot
=\frac{1}{2D}g_{kl}(\dot q^k+b^k)\partial_i b^l
-\tfrac12\,\partial_i\nabla^2\Pot,
\]
where we used that $g$ is constant (\Cref{def:manifold}), so $\partial_i g_{kl}
=0$.

\emph{Assembling.} Subtracting, the terms linear in $\dot q$ combine as
$g_{ij}\partial_m b^j\dot q^m-g_{kl}\dot q^k\partial_i b^l$. Because
$g_{ij}b^j=\partial_i\Pot$, we have $g_{ij}\partial_m b^j=\partial_m\partial_i
\Pot$, symmetric in $(i,m)$; the two first-order terms therefore cancel,
leaving
\[
\frac{1}{2D}g_{ij}\ddot q^j
-\frac{1}{2D}g_{kl}b^k\partial_i b^l
+\tfrac12\partial_i\nabla^2\Pot=0,
\]
which is \eqref{eq:EL-second-order} on writing $g_{kl}b^k\partial_i b^l=
\partial_i(\tfrac12 g_{kl}b^kb^l)$ (again using constancy of $g$).

\emph{Gradient limit.} The Lagrangian is a perfect square plus a potential
correction: $\Lag\ge-\tfrac12\nabla^2\Pot$ with equality iff $f^i=0$, i.e.\
$\dot q^i=-b^i=-g^{ij}\partial_j\Pot$. This zero-cost path is
\eqref{eq:gradflow}.
\end{proof}

\begin{corollary}[Component equations]\label{cor:components}
With $\Pot$ of \Cref{def:potential}, the gradient flow \eqref{eq:gradflow} reads
\begin{align}
\dot\Ord&=-\tfrac{1}{\mu_R}\!\left[(\Kc-K)\Ord+K\Ord^3-\alpha e^{-\sigma^2/2\pi^2}\right],\\
\dot{\sigma}^2&=-\tfrac{1}{\mu_\sigma}\!\left[\kB T-\frac{\kB T}{K\sigma^4}
+\frac{\alpha\Ord}{2\pi^2}e^{-\sigma^2/2\pi^2}\right],\\
\dot S_i&=-\partial_{S_i}V_{\mathrm{ent}}
=T_{\mathrm{eff}}\ln\!\frac{1-S_i}{S_i}-\lambda\,\partial_{S_i}(\text{wall}),
\quad i\in\{k,t,e\}.
\end{align}
\end{corollary}

\begin{proof}
Substitute the partials of \Cref{def:potential} into \eqref{eq:gradflow} with
$g^{-1}=\operatorname{diag}(\mu_R^{-1},\mu_\sigma^{-1},1,1,1)$.
\end{proof}

\subsection{The Kuramoto bifurcation and the corrected critical coupling}

\begin{theorem}[Synchronisation bifurcation]\label{thm:bifurcation}
For $K<\Kc$ the synchronisation potential $V_{\mathrm{sync}}(\cdot;K)$ has its
unique minimum at $\Ord=0$ (incoherence); for $K>\Kc$ a stable partially
synchronised minimum appears at $\Ord^*=\sqrt{1-\Kc/K}$. This is a continuous
(second-order) bifurcation, and it is the only genuine phase transition in the
model.
\end{theorem}

\begin{proof}
$\partial_\Ord V_{\mathrm{sync}}=(\Kc-K)\Ord+K\Ord^3=\Ord[(\Kc-K)+K\Ord^2]$,
vanishing at $\Ord=0$ and, for $K>\Kc$, at $\Ord^*=\sqrt{1-\Kc/K}$. The second
derivative $\partial^2_\Ord V_{\mathrm{sync}}=(\Kc-K)+3K\Ord^2$ is positive at
$\Ord=0$ iff $K<\Kc$ and is positive at $\Ord^*$ for $K>\Kc$; the minimum thus
moves continuously off zero as $K$ crosses $\Kc$, the signature of a
second-order transition.
\end{proof}

\begin{theorem}[Critical coupling]\label{thm:Kc}
For natural frequencies drawn from a unimodal symmetric density $g$ centred at
the mean, the mean-field critical coupling is
\[
\Kc=\frac{2}{\pi\,g(0)} .
\]
For a Gaussian law of standard deviation $\sigma_\omega$, $g(0)=1/(\sqrt{2\pi}
\,\sigma_\omega)$, giving
\[
\boxed{\;\Kc=\frac{2}{\pi}\sqrt{2\pi}\,\sigma_\omega=\frac{2\sqrt{2\pi}}{\pi}\,
\sigma_\omega\approx 1.596\,\sigma_\omega.\;}
\]
\end{theorem}

\begin{proof}
This is the classical Kuramoto mean-field result~\cite{kuramoto1975,strogatz2000,
acebron2005,dorfler2014}: linearising the self-consistency equation about
$\Ord=0^+$ gives the onset at $\Kc=2/[\pi g(0)]$. Substituting the Gaussian
$g(0)$ yields the boxed value.
\end{proof}

\begin{remark}[Correction of the preliminary coefficient]\label{rem:Kc-correction}
The preliminary version stated $\Kc=2\sigma_\omega/\pi\approx0.637\,\sigma_\omega$.
This is incorrect: it omits the factor $\sqrt{2\pi}$ from $g(0)$ for a Gaussian
and is too small by $\sqrt{2\pi}\approx2.507$. The correct value is
$\Kc=2/[\pi g(0)]$, equal to $1.596\,\sigma_\omega$ for a Gaussian. We confirm
the corrected law by direct simulation in \Cref{sec:external} (EXT02), where we
also show that $2\sigma_\omega/\pi$ fails badly for non-Gaussian frequency laws
and that the defensible general statement is $\Kc=2/[\pi g(0)]$, valid only for
unimodal symmetric $g$.
\end{remark}

\subsection{The five operational bands (not transitions)}

\begin{definition}[Operational coordination bands]\label{def:bands}
For descriptive convenience we label bands of the order parameter:
turbulent ($\Ord<0.3$), aperture-dominated ($0.3\le\Ord<0.5$), cascade
($0.5\le\Ord<0.8$), coherent ($0.8\le\Ord<0.95$), phase-locked
($\Ord\ge0.95$).
\end{definition}

\begin{proposition}[The band boundaries are not transitions]\label{prop:not-transitions}
The cut-offs $\Ord\in\{0.3,0.5,0.8,0.95\}$ are not critical points of $\Pot$:
for generic $K$, $\partial^2_\Ord\Pot$ does not change sign at these values, and
no thermodynamic singularity occurs there. They are classification cut-offs on
a continuous variable.
\end{proposition}

\begin{proof}
$\partial^2_\Ord V_{\mathrm{sync}}=(\Kc-K)+3K\Ord^2$ is smooth and strictly
increasing in $\Ord^2$ with at most one zero (at $\Ord=\sqrt{(K-\Kc)/3K}$ when
$K>\Kc$), which does not in general coincide with any of
$\{0.3,0.5,0.8,0.95\}$. The remaining terms of $\Pot$ are smooth in $\Ord$ and
introduce no sign change at the cut-offs. Hence no transition occurs at the band
boundaries.
\end{proof}

\begin{remark}
\Cref{prop:not-transitions} is the explicit retraction promised in
\Cref{sec:intro}. The bands remain useful as labels; they carry no claim of
criticality.
\end{remark}

%==========================================================
\part{Ensembles and Coordination}
%==========================================================

\section{Ensembles and the joint action}
\label{sec:ensemble}

\begin{definition}[Ensemble]\label{def:ensemble}
An \emph{ensemble} of $n$ agents is $E=(\agent_1,\dots,\agent_n)$ with a
coupling matrix $K_{ij}\ge0$ and natural-frequency vector
$\omega=(\omega_1,\dots,\omega_n)$ of standard deviation $\sigma_\omega$. Each
agent $\agent_i$ carries its own state $\qvec_i\in\M$ and internal order
parameter $\Ord_i$ with mean internal phase $\phi_i$.
\end{definition}

\begin{definition}[Ensemble order parameter]\label{def:ens-R}
The \emph{ensemble order parameter} is
\[
\Ord_{\mathrm{ens}}\,e^{i\psi_{\mathrm{ens}}}
=\frac1n\sum_{i=1}^n \Ord_i\,e^{i\phi_i},\qquad \Ord_{\mathrm{ens}}\in[0,1].
\]
\end{definition}

\begin{remark}[Hierarchical independence]
Each $\Ord_i$ measures intra-agent coherence; $\Ord_{\mathrm{ens}}$ measures
inter-agent coherence. The two are independent: an ensemble of internally
coherent agents ($\Ord_i\approx1$) can have $\Ord_{\mathrm{ens}}\approx0$ if
their phases $\phi_i$ are dispersed, and conversely.
\end{remark}

\begin{definition}[Joint action]\label{def:joint}
The ensemble dynamics follow the joint Onsager--Machlup action with Lagrangian
\[
\Lag_{\mathrm{ens}}=\sum_{i=1}^n \Lag^{(i)}+\Lag_{\mathrm{coup}},\qquad
\Lag_{\mathrm{coup}}=-\frac{1}{2n}\sum_{i,j}K_{ij}\cos(\phi_i-\phi_j),
\]
where $\Lag^{(i)}$ is the single-agent Lagrangian (\Cref{def:action}) and
$\Lag_{\mathrm{coup}}$ is the standard Kuramoto coupling.
\end{definition}

\begin{theorem}[Ensemble synchronisation]\label{thm:ens-sync}
The Euler--Lagrange equations of $\Lag_{\mathrm{ens}}$ reduce, in the gradient
limit, to per-agent gradient flow (\Cref{cor:components}) plus the Kuramoto
phase equation $\dot\phi_i=\omega_i-\frac1n\sum_j K_{ij}\sin(\phi_i-\phi_j)$;
the ensemble synchronises ($\Ord_{\mathrm{ens}}>0$ stable) for effective
coupling above the critical value of \Cref{thm:Kc} applied at the ensemble
level.
\end{theorem}

\begin{proof}
The single-agent part is \Cref{thm:EL}. For the coupling term,
$\partial\Lag_{\mathrm{coup}}/\partial\phi_i=\frac1n\sum_j K_{ij}\sin(\phi_i-
\phi_j)$, giving the Kuramoto phase equation; its mean-field bifurcation is
\Cref{thm:Kc} with $g$ the empirical frequency density.
\end{proof}

\begin{remark}[Equations of state are descriptive]
The preliminary version assigned each band a distinct ``equation of state''
$PV=N\kB T\,\mathcal{S}_{\mathrm{regime}}$. We retain these only as
\emph{descriptive summaries} of the structural factor in each band (e.g.\
$\mathcal{S}_{\mathrm{coh}}=1+\Ord^2/(1-\Ord^2)$ in the coherent band), and we
do not claim they are derived laws or that the bands they describe are
transitions. Their status is phenomenological.
\end{remark}

\section{Coordination on a common outcome}
\label{sec:coordination}

We give the minimal, defensible coordination statement: agents coordinate when
their order parameters lock, irrespective of internal representation. The
strong ``common-cell'' claim is tested empirically in \Cref{sec:external}
(EXT08).

\begin{definition}[Coordinated ensemble]\label{def:coordinated}
An ensemble is \emph{coordinated at level $\rho$} if $\Ord_{\mathrm{ens}}\ge\rho$.
\end{definition}

\begin{theorem}[Coordination is representation-independent]\label{thm:rep-indep}
$\Ord_{\mathrm{ens}}$ depends only on the agents' order parameters $\Ord_i$ and
phases $\phi_i$, not on their internal knowledge frameworks $\K_i$. Agents with
disjoint $\K_i$ can therefore attain any $\Ord_{\mathrm{ens}}$.
\end{theorem}

\begin{proof}
$\Ord_{\mathrm{ens}}$ is defined (\Cref{def:ens-R}) purely from $(\Ord_i,
\phi_i)$, which are scalar reductions of each agent's state; the frameworks
$\K_i$ do not enter. Hence disjointness of the $\K_i$ places no constraint on
$\Ord_{\mathrm{ens}}$.
\end{proof}

\begin{remark}
\Cref{thm:rep-indep} is the conservative core of the ``coordination without
shared representation'' idea. The stronger claim---that agents with disjoint
representations attain the \emph{same outcome label}---is an empirical question;
we test it directly on representation-disjoint classifiers in EXT08 and report
the (positive) result there.
\end{remark}

%==========================================================
\part{Partition Extinction}
%==========================================================

\section{An analogy-free theory of partition extinction}
\label{sec:extinction}

Reviewers asked for a definition of ``partition extinction'' that does not rely
on physical analogy or interpretive language, and for a clear separation
between \emph{functional} equivalence of agents and \emph{ontological}
indistinguishability of particles. We provide both. The physical instances
(superconductivity, superfluidity) are then \emph{examples} of the definition,
tested in \Cref{sec:external} (EXT07), not part of it.

\subsection{Partition operations and lag}

\begin{definition}[Partition operation and lag]\label{def:partition-op}
A \emph{partition operation} between two constituents $i,j$ is a procedure that
assigns them to distinct cells of \Cref{ax:resolution}. Its \emph{lag}
$\tauP^{(ij)}>0$ is the time the procedure requires. By finite resolution and
finite signalling speed, distinguishing states separated by $\Delta x$ requires
$\tauP^{(ij)}\ge\Delta x/c$ for a signalling speed $c$.
\end{definition}

\begin{definition}[Categorical distinguishability]\label{def:distinguishability}
Constituents $i,j$ are \emph{categorically distinguishable} if there exists a
partition operation assigning them to different cells; otherwise they are
\emph{categorically indistinguishable}. By \Cref{ax:resolution} this is a binary
property: there is no partial distinguishability below resolution $\delta$.
\end{definition}

\begin{definition}[Partition extinction]\label{def:extinction}
\emph{Partition extinction} between $i,j$ is the condition that $i,j$ are
categorically indistinguishable, so that no partition operation exists and the
lag $\tauP^{(ij)}$ is undefined; by convention we set $\tauP^{(ij)}=0$ to keep
transport quantities (in which $\tauP$ appears multiplicatively) well-defined.
\end{definition}

This definition uses only \Cref{ax:resolution} and \Cref{def:partition-op}; it
mentions no physical system.

\subsection{The extinction theorem}

\begin{theorem}[Partition extinction is discontinuous]\label{thm:extinction}
Let a control parameter $\theta$ (e.g.\ a coupling or an inverse temperature)
drive a pair from distinguishable to phase-locked, where phase-locking means
$\phi_i=\phi_j$ identically. Then the lag transitions discontinuously,
\[
\tauP^{(ij)}(\theta)=\begin{cases}\tauP^{(ij)}>0 & \text{distinguishable},\\
0 & \text{phase-locked},\end{cases}
\]
and any transport coefficient $\Theta=\frac1N\sum_{i,j}\tauP^{(ij)}G_{ij}$
(with bounded weights $G_{ij}$) attains $\Theta=0$ exactly in the phase-locked
phase.
\end{theorem}

\begin{proof}
\emph{Step 1 (binary).} By \Cref{def:distinguishability}, distinguishability is
binary: either a partition operation exists or it does not. \emph{Step 2
(phase-locking $\Rightarrow$ indistinguishable).} If $\phi_i=\phi_j$ identically,
the two constituents trace identical trajectories; by \Cref{ax:null} their
categorical state is fixed by phase-space location, so they occupy the same
categorical state and no operation distinguishes them. \emph{Step 3 (lag).} For
indistinguishable constituents no partition operation exists, so $\tauP^{(ij)}$
is undefined; the convention $\tauP^{(ij)}=0$ (\Cref{def:extinction}) is forced
by requiring $\Theta$ to remain defined. \emph{Step 4 (discontinuity).} Because
distinguishability is binary (Step 1), there is no value of $\theta$ at which
the pair is ``partly distinguishable''; the lag therefore jumps from a positive
value to $0$ without intermediate values. Substituting $\tauP^{(ij)}=0$ for all
locked pairs gives $\Theta=\frac1N\sum_{i,j}0\cdot G_{ij}=0$ exactly.
\end{proof}

\begin{remark}[``Exactly zero'' is a statement about the model]
The exactness in \Cref{thm:extinction} is a property of the idealised
categorical model, in which distinguishability is strictly binary. In any
physical realisation, residual distinguishing channels (finite-size effects,
disorder, noise, clock drift) keep a small positive $\tauP$; the model's exact
zero is the limit as those channels vanish. We are explicit about this: the
defensible empirical claim (EXT07) is that measured transport coefficients fall
\emph{below any detectable bound} (by $\ge 14$ orders of magnitude in
superconductors), which is consistent with, but does not prove, an exact zero.
\end{remark}

\subsection{Functional vs.\ ontological indistinguishability}

\begin{definition}[Categorical and physical observables]\label{def:observables}
A \emph{categorical observable} $\hat O_{\mathrm{cat}}$ is a function of the
partition coordinates $(n,\ell,m,s)$ alone. A \emph{physical observable}
$\hat O_{\mathrm{phys}}$ is a function of position, momentum, and energy alone.
\end{definition}

\begin{theorem}[Observable commutation]\label{thm:commutation}
Categorical and physical observables are functionally independent:
$[\hat O_{\mathrm{cat}},\hat O_{\mathrm{phys}}]=0$, in the sense that the value
of either is unchanged by conditioning on the other. Consequently, two agents
may be categorically indistinguishable (equal on all $\hat O_{\mathrm{cat}}$)
while remaining physically distinct (different on some $\hat O_{\mathrm{phys}}$),
and vice versa.
\end{theorem}

\begin{proof}
By \Cref{def:observables} the two observable classes are functions of disjoint
argument sets---partition coordinates versus phase-space coordinates---which by
\Cref{thm:commutation}'s hypothesis vary independently. A function of one set is
constant along variations of the other, so measuring one does not constrain the
other; their joint algebra is the tensor product and they commute. Equality on
all $\hat O_{\mathrm{cat}}$ (categorical indistinguishability) therefore implies
nothing about $\hat O_{\mathrm{phys}}$, establishing the final claim.
\end{proof}

\begin{remark}[Resolving the conflation]\label{rem:functional-ontological}
\Cref{thm:commutation} answers the reviewers' objection directly. When we say
synchronised agents become ``categorically indistinguishable,'' we mean
\emph{functionally} indistinguishable---equal on the categorical observables
that the coordination dynamics act upon. We do \emph{not} claim they lose
physical identity, computational independence, or distinct boundaries, as
bosons do under exchange. The preliminary version's Cooper-pair language
overstated this. The correct statement is: synchronised agents are equivalent
on $\hat O_{\mathrm{cat}}$ while remaining distinct on $\hat O_{\mathrm{phys}}$;
functional equivalence is not ontological identity.
\end{remark}

\subsection{Synchronisation as a structural correspondence}

\begin{proposition}[Coordination/extinction correspondence]\label{prop:correspondence}
In a phase-locked ensemble ($\Ord_{\mathrm{ens}}=1$, all $\phi_i$ equal), every
pair is categorically indistinguishable, so by \Cref{thm:extinction} the
inter-agent partition lag is zero and the coordination ``friction''
$\Theta_{\mathrm{coord}}=\frac1n\sum_{i,j}\tauP^{(ij)}G_{ij}$ vanishes in the
categorical model.
\end{proposition}

\begin{proof}
Phase-locking gives $\phi_i=\phi_j$ for all pairs; \Cref{thm:extinction} Step 2
gives indistinguishability and Step 3--4 give $\tauP^{(ij)}=0$, hence
$\Theta_{\mathrm{coord}}=0$.
\end{proof}

\begin{remark}[Scope of the correspondence]\label{rem:correspondence-scope}
\Cref{prop:correspondence} is a \emph{structural correspondence} between the
coordinated regime and partition extinction, not an identity between agent
coordination and physical condensation. Three caveats, retained from the
referee discussion: (i) the friction vanishes \emph{in the categorical model};
real coordinated systems retain communication delay, noise, and clock drift,
so the operational friction is small but positive (cf.\
\Cref{rem:functional-ontological}); (ii) we do \emph{not} assert that the
synchronisation transition is reversible in the hysteretic-physical sense---the
variational model is gradient and admits no claim about hysteresis, so we make
none; (iii) the correspondence transfers structure (binary distinguishability,
discontinuous lag), not physical mechanism (there is no exchange statistics for
agents). With these caveats the correspondence is defensible; without them it
overreaches, and we have removed the overreaching statements of the preliminary
version.
\end{remark}

%==========================================================
\part{External Validation, Scope, and Conclusion}
%==========================================================

\section{External validation against public data}
\label{sec:external}

\subsection{Why external, and what ``validation'' means here}

The preliminary version reported numerical checks that compared closed-form
predictions to the same closed forms; agreement ``to machine precision''
($\sim10^{-16}$) measured floating-point rounding of identities, not empirical
adequacy. We replace those with nine experiments against \emph{public data the
theory did not generate}. We retain the original closed-form checks only under
their correct name---\emph{internal consistency checks}---and do not count them
as validation.

Each experiment is fully reproducible; code and machine-readable results are
provided with the manuscript. We report each verdict honestly using three
labels: \textsc{confirmed} (data support the claim), \textsc{confirmed$^\ast$}
(supported but requiring an explicit scope or caveat), and
\textsc{correction} (data contradict the stated claim, forcing a change to the
theory). Two of the nine are corrections; we regard their presence as evidence
that the confirmations are not artefacts.

\subsection{The nine experiments}

\paragraph{EXT01 --- Shell capacity $C(n)=2n^2$ (atomic structure).}
The degeneracy law (\Cref{thm:capacity}) was compared to tabulated atomic shell
capacities, subshell sums $2(2\ell+1)$, and periodic-table period lengths.
\textsc{Confirmed$^\ast$}: $C(n)=2n^2$ matches the level degeneracy exactly for
$n=1$--$7$, and the Aufbau ($n+\ell$) order reproduces the real period lengths
$2,8,8,18,18,32,32$; however, pure $2n^2$ nesting reproduces only the He and Ne
closures, not Ar onward---shell \emph{closure} requires the filling rule, which
$2n^2$ alone does not supply. We state this limitation rather than claim closure
from degeneracy.

\paragraph{EXT02 --- Kuramoto critical coupling (direct simulation).}
We integrated the finite-$N$ Kuramoto dynamics ($N=400$) for four
frequency laws and located the synchronisation onset by extrapolating the
$\Ord^2$ takeoff. \textsc{Correction}: the preliminary $\Kc=2\sigma_\omega/\pi$
is wrong by $\sqrt{2\pi}\approx2.51$; the simulated onset is consistent with the
corrected $\Kc=2/[\pi g(0)]$ ($=1.596\,\sigma_\omega$ for a Gaussian, the value
adopted in \Cref{thm:Kc}). For the Lorentzian law---where the mean-field result
is exact---the prediction $\Kc=2\gamma$ matched the simulation, while
$2\sigma_\omega/\pi$ erred by a factor exceeding twenty; for uniform and bimodal
laws $2\sigma_\omega/\pi$ also failed. The corrected formula is adopted
throughout and scoped to unimodal symmetric $g$.

\paragraph{EXT03 --- Wiedemann--Franz constant (eleven metals).}
The framework reproduces the universality of the Lorenz ratio
$L=\kappa/(\sigma T)$. \textsc{Confirmed}: the Sommerfeld value
$L_0=(\pi^2/3)(\kB/e)^2=2.44\times10^{-8}\,\mathrm{W\,\Omega\,K^{-2}}$ was
compared to measured $L$ for eleven metals at $0$ and $100^\circ$C
\cite{kittel2005}; the measured mean is $2.51\times10^{-8}$ with $20$ of $22$
values within $15\%$. We note explicitly that this \emph{reproduces} standard
physics and is not a novel prediction; the framework is consistent with WF, it
does not improve on it.

\paragraph{EXT04 --- Sleep-stage order parameter (real EEG).}
The companion claim that sleep stages order as $\Ord_{N3}>\Ord_W>\Ord_{N1}>
\Ord_{N2}>\Ord_{\mathrm{REM}}$ was tested on a polysomnography recording from
PhysioNet Sleep-EDF~\cite{kemp2000,goldberger2000} by computing the
cross-channel Kuramoto order parameter per $30$\,s epoch.
\textsc{Correction}: the observed ordering is $\Ord_W>\Ord_{\mathrm{REM}}>
\Ord_{N1}>\Ord_{N2}>\Ord_{N3}$---the \emph{reverse} of the claim, with N3
lowest, not highest. A caveat applies (the recording has only two bipolar EEG
derivations, so cross-channel $\Ord$ measures fronto-parietal coherence rather
than within-population synchrony), but on this evidence the stage-ordering claim
is unsupported. We therefore \emph{withdraw} the sleep--regime ordering as a
claim of the framework and mention sleep only as a possible application
requiring proper multichannel testing.

\paragraph{EXT05 --- Enzyme efficiency vs.\ substrate complexity.}
The claim that catalytic efficiency anti-correlates with transformation
complexity was tested using a framework-independent complexity proxy (substrate
heavy-atom count) against literature $k_{\mathrm{cat}}/K_M$~\cite{fersht1999,
bionumbers2015}. \textsc{Confirmed$^\ast$}: $\log(k_{\mathrm{cat}}/K_M)$
anti-correlates with substrate heavy-atom count, Pearson $r=-0.67$ ($p=0.05$),
matching the claimed $\approx-0.57$ in sign and magnitude; but the trend is weak
($r^2=0.45$, Spearman $-0.40$) and driven by the extremes. The defensible claim
is a qualitative anti-correlation, not a quantitative law.

\paragraph{EXT06 --- Antidepressant response convergence.}
The claim that mechanistically distinct drugs converge on similar efficacy was
tested against the Cipriani et al.\ network meta-analysis~\cite{cipriani2018}:
response odds ratios for eleven antidepressants spanning SSRI/SNRI/TCA/atypical
classes, versus their receptor-binding breadth. \textsc{Confirmed$^\ast$}:
response ORs cluster tightly (coefficient of variation $0.10$) and binding
breadth does not predict response (Pearson $r=0.39$, $p=0.23$, n.s.), supporting
breadth-independent convergence. We note this is correlational and consistent
with several explanations; the aperture/regime mechanism is one interpretation,
not a demonstrated cause.

\paragraph{EXT07 --- Partition extinction (superconductors).}
The claim of exact, discontinuous transport vanishing (\Cref{thm:extinction})
was tested against measured superconducting resistivity bounds and critical
temperatures. \textsc{Confirmed$^\ast$}: measured superconducting resistivity
upper bounds lie $\ge14$ orders of magnitude below the normal state
(operationally exact zero) at a sharp $\Tc$ across ten superconductors, plus the
He-4 $\lambda$-point where viscosity vanishes. As in EXT03, this reproduces the
BCS/Landau phenomenology and does not predict $\Tc$ from first principles; the
framework is consistent with the exact-vanishing/sharp-transition phenomenology,
as an interpretation alongside the standard theory.

\paragraph{EXT08 --- Common-cell convergence (disjoint representations).}
The central coordination claim---agents with disjoint representations attain the
same outcome---was tested on the scikit-learn digits benchmark
\cite{pedregosa2011} using four classifiers trained on \emph{mutually disjoint}
pixel quadrants (no shared feature). \textsc{Confirmed}: the four agents agree
on the predicted label at rate $0.51$ versus chance $0.10$; on the $29\%$ of
inputs where all four agree, they are correct $100\%$ of the time. This is a
genuine instance of coordination on the outcome (the label) without shared
representation, supporting \Cref{thm:rep-indep} and the common-cell idea.

\paragraph{EXT09 --- Catalytic composition (real classifiers).}
The composition law---combining independent agents lowers the error floor below
the best individual---was tested by majority-voting the EXT08 agents.
\textsc{Confirmed$^\ast$}: ensemble error ($0.16$) is below the best individual
error ($0.26$), so composition does lower the floor. However, the \emph{exact}
multiplicative prediction (which assumes independence) lies far below the
measured value because the agents are not independent (mean pairwise
error-correlation $0.15$; they fail on the same hard digits). We therefore
state the composition law as an \emph{independence-limit bound}, not an
equality.

\subsection{Summary of verdicts}

\begin{table}[h!]\centering
\small
\begin{tabular}{llp{0.36\columnwidth}}
\toprule
\textbf{Exp.} & \textbf{Verdict} & \textbf{One-line outcome}\\\midrule
EXT01 & confirmed$^\ast$ & $2n^2$ exact as degeneracy; closure needs Aufbau\\
EXT02 & \textbf{correction} & $\Kc=2\sigma/\pi$ wrong by $\sqrt{2\pi}$; use $2/[\pi g(0)]$\\
EXT03 & confirmed & $L_0$ matches 11 metals (20/22 within 15\%)\\
EXT04 & \textbf{correction} & sleep order inverts on real EEG; claim withdrawn\\
EXT05 & confirmed$^\ast$ & enzyme anti-correlation $r=-0.67$, weak\\
EXT06 & confirmed$^\ast$ & response convergence CV $=0.10$, breadth n.s.\\
EXT07 & confirmed$^\ast$ & resistivity $\downarrow\ge10^{14}$; reproduces BCS\\
EXT08 & confirmed & disjoint agents 100\% correct when agreeing\\
EXT09 & confirmed$^\ast$ & ensemble beats best; independence violated\\
\bottomrule
\end{tabular}
\caption{External-validation verdicts. Two outright confirmations, five
confirmations requiring explicit scope, and two corrections forced by the data
(EXT02, EXT04). Full code and JSON results accompany the manuscript.}
\label{tab:external}
\end{table}

\section{Scope and limitations}
\label{sec:scope}

We delimit what the framework does and does not establish.

\begin{enumerate}[leftmargin=*]
\item \emph{The axioms are general, hence weakly discriminating.}
\Cref{ax:bounded,ax:null,ax:resolution} hold for almost all physical and
computational systems. Their function is to fix the state space and counting
structure; discriminating content lies in the dynamics and in the external
tests, not in the axioms.

\item \emph{S-entropy coordinates are distributional, not content-laden.} They
measure the informational geometry of an agent's state
(\Cref{rem:s-content}); they do not encode ecological salience or semantic
content. The framework therefore does not model what specific information a
biological agent attends to, only the shape of its state distribution.

\item \emph{The model is deterministic-overdamped and mean-field.} Real neural
and collective systems are strongly stochastic and multiscale, with fluctuating
network structure. We capture fluctuations only through the Onsager--Machlup
noise term and do not model metastability or time-varying connectivity. Claims
about such systems are therefore at the level of qualitative structure, not
quantitative dynamics.

\item \emph{The band boundaries are operational.} As proved in
\Cref{prop:not-transitions}, $\{0.3,0.5,0.8,0.95\}$ are classification cut-offs,
not transitions; only the Kuramoto bifurcation is a genuine transition.

\item \emph{Partition extinction is exact only in the categorical model.} The
operational friction of real coordinated systems is small but positive
(\Cref{rem:correspondence-scope}); we make no hysteresis or reversibility
claim.

\item \emph{Physical instances reproduce, they do not predict.} For
Wiedemann--Franz (EXT03) and superconductivity (EXT07), the framework recovers
standard results; it does not yield new quantitative predictions there, and we
do not claim it does.

\item \emph{Withdrawn claim.} The sleep-stage ordering (EXT04) is not supported
by real EEG and is withdrawn as a claim of the framework.

\item \emph{Composition requires independence.} The multiplicative composition
law (EXT09) holds as an independence-limit bound; real agents that share a task
violate independence and compose less favourably.
\end{enumerate}

\section{Conclusion}
\label{sec:conclusion}

We have rebuilt the coordination calculus as a self-contained theory. From three
axioms we derived the S-entropy manifold (via the triple-equivalence identity),
the partition potential, and the agent action, and we gave the Euler--Lagrange
equations explicitly rather than by assertion. We corrected the critical-coupling
coefficient to the standard $\Kc=2/[\pi g(0)]$, demoted the five regime
boundaries from phase transitions to operational cut-offs (the only genuine
transition being the Kuramoto bifurcation), and replaced the analogy-laden
account of partition extinction with a formal definition and an
observable-commutation theorem that separates functional from ontological
indistinguishability. Most importantly, we replaced self-referential numerical
checks with nine external tests against public data; two forced corrections to
the theory, and the remainder confirm it within stated scope. What the framework
defensibly offers is a variational language in which coordination is
order-parameter synchronisation, the deepest coordinated regime corresponds
structurally to partition extinction, and agents with disjoint representations
can nonetheless converge on a common outcome---this last now demonstrated on
real data rather than asserted.

\bibliographystyle{plain}
\bibliography{references}

\end{document}
